\section{USAMO 2006}
        \begin{enumerate}
        	\item (\textit{USAMO 2006}) Let $p$ be a prime number and let $s$ be an integer with $0<s<p$. Prove that there exist integers $m$ and $n$ with $0<m<n<p$ and

 \[\left\lbrace\frac{sm}{p}\right\rbrace<\left\lbrace\frac{sn}{p}\right\rbrace<\frac{s}{p}\]
if and only if $s$ is not a divisor of $p-1$.


 Note: For $x$ a real number, let $\lfloor x\rfloor$ denote the greatest integer less than or equal to $x$, and let $\lbrace x\rbrace=x-\lfloor x\rfloor$ denote the fractional part of $x$.


 

	\item (\textit{USAMO 2006}) For a given positive integer $k$ find, in terms of $k$, the minimum value of $N$ for which there is a set of $2k+1$ distinct positive integers that has sum greater than $N$ but every subset of size $k$ has sum at most $N/2$.


 

	\item (\textit{USAMO 2006}) For integral $m$, let $p(m)$ be the greatest prime divisor of $m$. By convention, we set $p(\pm1)=1$ and $p(0)=\infty$. Find all polynomials $f$ with integer coefficients such that the sequence $\lbrace p(f(n^2))-2n\rbrace_{n\ge0}$ is bounded above. (In particular, this requires $f(n^2)\neq0$ for $n\ge0$.)


 

	\item (\textit{USAMO 2006}) Find all positive integers $n$ for which there exists an integer $k\ge 2$ and positive rationals $a_1,a_2,\ldots a_k$ satisfying $a_1+a_2+\ldots+a_k=a_1\cdot a_2\cdots a_k=n$.


 

	\item (\textit{USAMO 2006}) A mathematical frog jumps along the number line. The frog starts at 1, and jumps according to the following rule: if the frog is at integer $n$, then it can jump either to $n+1$ or to $n+2^{m_n+1}$ where $2^{m_n}$ is the largest power of 2 that is a factor of $n$. Show that if $k\ge2$ is a positive integer and $i$ is a nonnegative integer, then the minimum number of jumps needed to reach $2^ik$ is greater than the minimum number of jumps needed to reach $2^i$.


 

	\item (\textit{USAMO 2006}) Let $ABCD$ be a quadrilateral, and let $E$ and $F$ be points on sides $AD$ and $BC$, respectively, such that $AE/ED=BF/FC$.  Ray $FE$ meets rays $BA$ and $CD$ at $S$ and $T$ respectively. Prove that the circumcircles of triangles $SAE$, $SBF$, $TCF$, and $TDE$ pass through a common point.


 

\end{enumerate}